% ============================================================
%  SCHOLA DOMESTICA — Magister Mathematica
%  Level 1 | Week 9 | Day 4
%  Topic: Addition Facts Sums to 6 — Vertical Form & Missing Addend
% ============================================================
\documentclass[12pt, letterpaper]{article}

\usepackage[top=1in, bottom=1in, left=1in, right=1in]{geometry}
\usepackage[T1]{fontenc}
\usepackage[utf8]{inputenc}
\usepackage{lmodern}
\usepackage{microtype}
\usepackage{amsmath}
\usepackage{xcolor}
\usepackage{tikz}
\usepackage{array}
\usepackage{enumitem}
\usepackage{multicol}
\usepackage{mdframed}
\usepackage{fancyhdr}

\definecolor{scholared}{RGB}{160,30,30}
\definecolor{scholagold}{RGB}{180,140,40}
\definecolor{scholacream}{RGB}{253,248,235}
\definecolor{scholagreen}{RGB}{50,110,60}
\definecolor{scholablue}{RGB}{40,80,140}

\pagestyle{fancy}
\fancyhf{}
\renewcommand{\headrulewidth}{0.6pt}
\fancyhead[L]{\small\textcolor{scholared}{\textbf{Schola Domestica — Magister Mathematica}}}
\fancyhead[R]{\small\textcolor{scholared}{Level 1 · Week 9 · Day 4}}
\fancyfoot[C]{\small\textcolor{gray}{— \thepage\ —}}

\newcommand{\answerline}{\underline{\hspace{2.5cm}}}
\newcommand{\biganswerline}{\underline{\hspace{4cm}}}
\newcommand{\numbox}{\tikz{\draw[scholablue,thick](0,0)rectangle(1.2cm,0.85cm);}}

% vertical addition display helper
\newcommand{\vertadd}[3]{%
  \begin{array}{r}
    #1 \\
  + #2 \\
  \hline
    #3
  \end{array}
}

\newmdenv[
  backgroundcolor=scholacream,
  linecolor=scholared,
  linewidth=1.5pt,
  roundcorner=6pt,
  innertopmargin=8pt,
  innerbottommargin=8pt,
  innerleftmargin=10pt,
  innerrightmargin=10pt
]{lessonbox}

\newmdenv[
  backgroundcolor={rgb,255:red,235;green,245;blue,235},
  linecolor=scholagreen,
  linewidth=1.2pt,
  roundcorner=4pt,
  innertopmargin=6pt,
  innerbottommargin=6pt,
  innerleftmargin=8pt,
  innerrightmargin=8pt
]{enrichbox}

\begin{document}

\begin{center}
  {\LARGE\textbf{\textcolor{scholared}{Sums to Six:}}}\\[4pt]
  {\Large\textcolor{scholagold}{\textit{Standing Up and Finding What's Missing}}}\\[6pt]
  {\small Name: \biganswerline \hspace{1cm} Date: \answerline}
\end{center}

\vspace{4pt}
\noindent\textcolor{scholared}{\rule{\linewidth}{1pt}}
\vspace{4pt}

% IMAGE PROMPT (Beatrix Potter watercolour style):
% A little mouse in a blue coat stacks acorns into a tower on a wooden table, counting on its paws.
% Cosy kitchen background with a candle. Warm amber palette, white paper background, no text.

\begin{lessonbox}
\textbf{Addition can be written two ways — they mean the same thing!}\\[6pt]
\noindent
\begin{tabular}{@{}p{5cm} p{5cm}@{}}
\textbf{Across (horizontal):} & \textbf{Down (vertical):} \\[4pt]
$3 + 2 = 5$ &
$\displaystyle\vertadd{3}{2}{5}$
\end{tabular}

\medskip
\noindent When you write addition \emph{vertically}, line up the digits and add down.
\end{lessonbox}

\bigskip

% ===== PART I: VERTICAL FORM =====
\noindent\textcolor{scholared}{\large\textbf{Part I — Solve the Vertical Problems}} \hfill \textit{(Problems 1–8)}

\medskip
\noindent\textit{Add the numbers. Write the sum on the line below.}

\bigskip

\noindent
\begin{tabular}{@{}p{1.8cm} p{1.8cm} p{1.8cm} p{1.8cm} p{1.8cm} p{1.8cm} p{1.8cm} p{1.8cm}@{}}

% 1: 2+1=3
\textbf{1.}
$\begin{array}{r} 2 \\ +1 \\ \hline \phantom{0} \end{array}$
&
% 2: 3+0=3
\textbf{2.}
$\begin{array}{r} 3 \\ +0 \\ \hline \phantom{0} \end{array}$
&
% 3: 1+3=4
\textbf{3.}
$\begin{array}{r} 1 \\ +3 \\ \hline \phantom{0} \end{array}$
&
% 4: 4+2=6
\textbf{4.}
$\begin{array}{r} 4 \\ +2 \\ \hline \phantom{0} \end{array}$
&
% 5: 2+3=5
\textbf{5.}
$\begin{array}{r} 2 \\ +3 \\ \hline \phantom{0} \end{array}$
&
% 6: 5+1=6
\textbf{6.}
$\begin{array}{r} 5 \\ +1 \\ \hline \phantom{0} \end{array}$
&
% 7: 0+0=0
\textbf{7.}
$\begin{array}{r} 0 \\ +0 \\ \hline \phantom{0} \end{array}$
&
% 8: 3+3=6
\textbf{8.}
$\begin{array}{r} 3 \\ +3 \\ \hline \phantom{0} \end{array}$

\end{tabular}

\bigskip\bigskip

% ===== PART II: MISSING ADDEND =====
\noindent\textcolor{scholared}{\large\textbf{Part II — What Is Missing?}} \hfill \textit{(Problems 9–12)}

\medskip

\begin{lessonbox}
\textbf{Missing addend:} One of the parts is hidden. Think: \textit{``What do I add to reach the whole?''}\\[4pt]
\hspace{2em} $3 + \boldsymbol{?} = 5$ \quad Think: $3 + 2 = 5$, so the missing part is \textbf{2}.
\end{lessonbox}

\bigskip

\noindent\textit{Find the missing addend. Write it in the box.}

\bigskip

\noindent
\begin{tabular}{@{}p{7.0cm} p{7.0cm}@{}}
\textbf{9.}\quad $2 +$ \numbox $= 4$ &
\textbf{10.}\quad $1 +$ \numbox $= 4$ \\[16pt]
\textbf{11.}\quad $3 +$ \numbox $= 6$ &
\textbf{12.}\quad \numbox $+ 2 = 5$
\end{tabular}

\bigskip\bigskip

% ===== PART III: WORD PROBLEMS =====
\noindent\textcolor{scholared}{\large\textbf{Part III — Story Problems}} \hfill \textit{(Problems 13–16)}

\medskip

\noindent\textbf{13.} A nest has \textbf{2} blue eggs and \textbf{3} speckled eggs. How many eggs in all?

\medskip
\noindent\hspace{2em}$\numbox + \numbox = \numbox$ eggs

\bigskip

\noindent\textbf{14.} Tom had \textbf{4} buttons. He found some more. Now he has \textbf{6}. How many did he find?

\medskip
\noindent\hspace{2em}$4 +$ \numbox $= 6$\quad He found \numbox\ buttons.

\bigskip

\noindent
\begin{tabular}{@{}p{7.0cm} p{7.0cm}@{}}
\textbf{15.} $5 + 0 =$ \numbox &
\textbf{16.} $0 + 6 =$ \numbox
\end{tabular}

\bigskip

\begin{enrichbox}
\textbf{\textcolor{scholagreen}{Challenge: Commutative Property}}\\[3pt]
Do these two facts give the same answer?
\quad $2 + 3 =$ \numbox \quad and \quad $3 + 2 =$ \numbox\\[3pt]
\textit{When we switch the parts, the whole stays the same. This is called the \textbf{turn-around rule}.}
\end{enrichbox}

\vspace{0.3cm}
\noindent\textcolor{gray}{\small\textit{Ray's Primary Arithmetic — vertical notation and missing addend.}}

\end{document}
